\scanpage{207}
\mspage{81}
\deleted{Mais} D'autre part, je dis que le syst. inductif des $f_i$ dans $\underline{\mathrm{Fl}}\,M$ est \uline{bien adapté à} $\unsure{\pi|R}$. On sait déjà que $I$ est \deleted{\unsure{a}} \unsure{bien}-adapté à $\pi$, le \deleted{\unsure{sup}} sup \inserted{dans $I$} d'une partie filtrante $J$ de cardinal $\leq\pi$ étant le sup dans l'ens. ordonné $\mathrm{Ss}_{\deleted{\unsure{\pi}}}(B)$ (lequel sup se trouve dans $\mathrm{Ss}_\pi(B)$). Il est essentiel, pour notre argument, de savoir que si $i_J=\sup J$, alors non seulement $B_{i_J}=\underset{i\in J}{\unsure{\underline{\lim}|\underrightarrow{\lim}}} B_i$, mais aussi $A_{i_J}=\underset{i\in J}{\underrightarrow{\lim}} A_i$, de sorte que \margin{$\Vert$}le système inductif des $f_i$ est $\pi$-adapté dans $\underline{\mathrm{Fl}}\,M$. Il semble qu'il faille donc une condition d'exactitude sur $M$ voisine de e, savoir

\textcircled{e'} les foncteurs $\underline{\lim}$ filtrantes commutent au chgt de base \deleted{\unsure{(dans}} \deleted{\unsure{(par un}}

(il \unsure{le|en} suffit ici dans le cas très particulier d'un chgt. de base par un mono \inserted{$A\to B$}, et d'une limite inductive filtr. de sous-objets de $B$ pour des morphismes \unsure{d'inclusions|d'inclusion} comme transitions. Notons que \textcircled{e}, \textcircled{e'} se trouvent généralement comme conséquences de la condition

(e$_0$) les $\underline{\lim}$ filtrantes dans $I$ sont exactes à gauche.
% archive sheet number: 206
